\chapter{Conclusion and Future Work}
\section{Answers to the Research Questions}
\textbf{RQ1:} Training-label corruption reduces aggregate performance and
its effect depends strongly on dataset, model, and transition structure.
The mean clean-reference accuracy is \CleanAccuracy{}, compared with
\NoisyAccuracy{} for noisy training. High-rate cyclic corruption is
particularly damaging on Digits and Dry Bean. Fixed-default HGB is more
sensitive than the other evaluated families on the dataset-level means,
but this is not a conclusion about optimally tuned boosting.

\textbf{RQ2:} Out-of-fold confidence is the most effective of the four
evaluated detectors on average, with detection F1 of 0.8072. Combined
scoring follows at 0.7504. General uncertainty, measured by the top-two
probability margin, is a poor proxy for injected label errors in this design.
The shared out-of-fold matrix and common selection budget make the component
comparison directly interpretable.

\textbf{RQ3:} Confidence-based removal improves average test accuracy by
\ConfidenceRemoveGain{} percentage points and combined removal by
\CombinedRemoveGain{} points. Neither closes the full gap to clean training.
Relabeling helps on average for three detectors, but is not uniformly better
than removal. High detection quality does not guarantee a downstream accuracy
gain, as the saved Dry Bean example demonstrates. Margin-only cleaning is
harmful on all four dataset-level averages.

\section{Contribution Supported by the Evidence}
The project delivers a reproducible benchmark and a transparent evaluation of
a simple combined-score hypothesis. The hypothesis does not succeed against
confidence alone on the aggregate comparison. The binary analysis explains
why the combination often cannot change confidence's ranking; the multiclass
experiments show that changing the ranking does not necessarily improve it.
These results favor starting with a simple confidence baseline before adding
uncertainty terms, within the controlled setting tested here.

The release includes an auditable matrix, generated tables and figures,
illustrative saved models, tests, pinned dependencies, and an immutable
Git-LFS archive. This supports verification and extension of the experiment,
not a claim of universal method superiority. With four datasets, no exact
two-sided comparison can reach significance at 0.05. Conclusions remain
descriptive and exploratory.

\section{Priorities for Future Work}
The first priority is to remove the oracle-budget assumption. A future study
should choose thresholds or removal fractions using only training information
or a properly allocated validation set, then evaluate the complete selection
procedure on untouched test data. This would test whether the current gains
survive uncertainty about the actual noise prevalence.

Second, the comparison should add established label-quality methods such as
confident learning, alternative cyclic maps, and feature-dependent corruptions.
Real annotation-error datasets would distinguish injected-error recovery from
genuine data-quality improvement. Additional mixed-type tabular datasets
would broaden the empirical population and improve the resolution of
dataset-level inference.

Third, model and score choices should be evaluated under a prespecified tuning
protocol. Probability calibration, class-aware selection budgets, and
confidence-gated relabeling are reasonable hypotheses, but should not be
selected from the current test outcomes and presented as already validated.
Randomized tie handling and separately seeded random processes are additional
sensitivity checks. Finally, controlled runtime and memory measurements
would turn the qualitative term ``lightweight'' into a measured cost comparison.
